% Part: turing-machines
% Chapter: undecidability
% Section: universal-tm

\documentclass[../../../include/open-logic-section]{subfiles}

\begin{document}

\olfileid{tur}{und}{uni}
\olsection{آلات تورنغ الشاملة}

قد تقدم في \olref[enu]{sec} أن لكل آلة تورنغ وصفًا بمتتالية
منتهية من الأعداد الصحيحة، تتضمن ترميز الحالات والأبجدية والحالة
الابتدائية والتعليمات. وتقدم أن مجموعة هذه الأوصاف !!{enumerable}.
ولكونها !!{denumerable} توجد دالة !!a{surjective} من~$\Nat$ إليها؛
ويمكن تحصيل دالة !!a{surjective} بهذه الصفة بطريقة كانتور المتعرجة،
مثلًا. فبذلك تتعدد أوصاف الآلات على نظام معين. وإذا ثبتنا هذا
التعداد، صح أن نقول: الآلة الأولى، والثانية،~\dots، والآلة ذات
الرتبة~$\OLId{latin-ordinary}{e}$. وتسمى هذه الأعداد \emph{فهارس}.

\begin{defn}
لتكن $\OLId{latin-looped}{M}$ الآلة ذات الرتبة~$\OLId{latin-ordinary}{e}$ في التعداد المثبت؛ عندئذ يسمى
$\OLId{latin-ordinary}{e}$~\emph{فهرسًا} للآلة~$\OLId{latin-looped}{M}$. ونكتب $\OLId{latin-looped}{M}_\OLId{latin-ordinary}{e}$ للدلالة على الآلة ذات
الرتبة~$\OLId{latin-ordinary}{e}$.
\end{defn}

ولا يلزم أن يكون للآلة فهرس واحد؛ فقد يختلف وصفان للآلة~$\OLId{latin-looped}{M}$
في ترتيب التعليمات وحده، فيقع كل وصف في موضع غير موضع الآخر،
ويختلف فهرساهما.

والأمر الذي نعتمد عليه أن هذا التعداد يمكن ترتيبه بحيث يُحسب وصف
$\OLId{latin-looped}{M}$ فعليًا من فهرسها، ويُحسب فهرس للآلة~$\OLId{latin-looped}{M}$ فعليًا من وصفها.
فتقتضي أطروحة تشرتش--تورنغ وجود آلة تورنغ تستخرج وصف الآلة ذات
الفهرس~$\OLId{latin-ordinary}{e}$ وتخرجه على شريطها. ويكون ذلك الوصف متتالية من كتل
علامات~$\TMstroke$، تمثل الأعداد الصحيحة الموجبة التي يتألف منها
وصف~$\OLId{latin-looped}{M}_\OLId{latin-ordinary}{e}$.

فأي الدوال على فهارس آلات تورنغ تقبل هي نفسها الحساب بهذه الآلات؟
وأي خواص تلك الفهارس يمكن تقريرها؟ من أمثلة ما يقبل الحساب الدالة
التي تعطي، من الفهرس~$\OLId{latin-ordinary}{e}$، عدد حالات الآلة ذات الفهرس~$\OLId{latin-ordinary}{e}$.
وبيان عمل آلة تحسبها أنها تبدأ بكتلة واحدة من $\OLId{latin-ordinary}{e}$ علامات
$\TMstroke$، فتستخرج من $\OLId{latin-ordinary}{e}$ الوصف الذي يرمزه. ويصير هذا الوصف
متتالية من كتل علامات~$\TMstroke$ على الشريط. ولما كان أول
!!{element} من المتتالية عدد الحالات، لم يبق عليها إلا محو ما سوى
الكتلة الأولى من علامات~$\TMstroke$، ثم التوقف.

ومن أهم ما يترتب على هذا التمثيل المبرهنة الآتية:

\begin{thm}\ollabel{thm:universal-tm}
توجد \emph{آلة تورنغ شاملة}~$\OLId{latin-looped}{U}$، إذا ابتدأت بالدخل $\tuple{\OLId{latin-ordinary}{e},\OLId{latin-ordinary}{n}}$،
كان شأنها ما يأتي:
  \begin{enumerate}
    \item تتوقف إذا وفقط إذا توقفت $\OLId{latin-looped}{M}_\OLId{latin-ordinary}{e}$ عند الدخل~$\OLId{latin-ordinary}{n}$؛
    \item وإذا توقفت $\OLId{latin-looped}{M}_\OLId{latin-ordinary}{e}$ بخرج $\OLId{latin-ordinary}{m}$، أخرجت~$\OLId{latin-looped}{U}$ العدد نفسه ثم توقفت.
  \end{enumerate}
  فالدالة التي تحسبها $\OLId{latin-looped}{U}$ هي $\OLId{latin-ordinary}{f}\colon\Nat\times\Nat\pto\Nat$،
  حيث~$\OLId{latin-ordinary}{f}(\OLId{latin-ordinary}{e},\OLId{latin-ordinary}{n})=\OLId{latin-ordinary}{m}$ متى توقفت $\OLId{latin-looped}{M}_\OLId{latin-ordinary}{e}$ عند الدخل~$\OLId{latin-ordinary}{n}$ بخرج~$\OLId{latin-ordinary}{m}$،
  وتكون غير معرّفة فيما عدا ذلك.
\end{thm}

\begin{proof}
  أما تفصيل إنشاء~$\OLId{latin-looped}{U}$ فشديد التعقيد حتى يكاد يتعذر عرضه هنا؛
  فنكتفي ببيان طريقة عملها، ثم نستند إلى أطروحة تشرتش--تورنغ.
  يكون على شريط~$\OLId{latin-looped}{U}$ أولًا كتلة من $\OLId{latin-ordinary}{e}$ علامات $\TMstroke$،
  وبعدها كتلة من $\OLId{latin-ordinary}{n}$ علامات~$\TMstroke$. فتفك ترميز الفهرس~$\OLId{latin-ordinary}{e}$
  وتضع الوصف إلى يمين الدخل~$\OLId{latin-ordinary}{n}$. وهذا الوصف قائمة أعداد، أي كتل
  علامات~$\TMstroke$ تفصلها رموز~$\TMblank$، تتضمن تعليمات~$\OLId{latin-looped}{M}_\OLId{latin-ordinary}{e}$.
  ثم تكتب $\OLId{latin-looped}{U}$ عدد الحالة الابتدائية للآلة~$\OLId{latin-looped}{M}_\OLId{latin-ordinary}{e}$ والعدد~$\OLMathNumeral{1}$
  إلى يمين وصف~$\OLId{latin-looped}{M}_\OLId{latin-ordinary}{e}$؛ وتمثلهما أيضًا تمثيلًا أحاديًا بكتل
  علامات~$\TMstroke$. وبعد ذلك تنسخ الدخل، وهو كتلة من $\OLId{latin-ordinary}{n}$
  علامات~$\TMstroke$، إلى اليمين، جاعلةً بدل كل علامة~$\TMstroke$
  كتلة من ثلاث علامات~$\TMstroke$. ذلك أن رمز $\TMstroke$
  عدده~$\OLMathNumeral{3}$، ورمز $\TMendtape$ عدده~$\OLMathNumeral{1}$، ورمز $\TMblank$
  عدده~$\OLMathNumeral{2}$. وتضع على يسار متتالية هذه الكتل، التي تفصلها
  رموز~$\TMblank$ في شريط $\OLId{latin-looped}{U}$، علامة~$\TMstroke$ واحدة،
  وهي ترميز~$\TMendtape$.

  فصار شريط~$\OLId{latin-looped}{U}$ يحمل الفهرس~$\OLId{latin-ordinary}{e}$، والعدد~$\OLId{latin-ordinary}{n}$، ورمز الحالة
  الابتدائية الذي يمثل «الحالة الحالية»، وعدد موضع الرأس الابتدائي~$\OLMathNumeral{1}$
  الذي يمثل «موضع الرأس الحالي»، ومحتويات «الشريط» الابتدائية:
  متتالية كتل من علامات~$\TMstroke$ تمثل أعداد رموز~$\OLId{latin-looped}{M}_\OLId{latin-ordinary}{e}$،
  ونسميها «الرموز»، وتفصل بينها علامات~$\TMblank$.

  ولتحاكي ما تفعله $\OLId{latin-looped}{M}_\OLId{latin-ordinary}{e}$ عند الدخل~$\OLId{latin-ordinary}{n}$، تجري على هذه الطريقة:
  \begin{enumerate}
    \item تطلب العدد~$\OLId{latin-ordinary}{k}$ المسجل ل«موضع الرأس الحالي»؛ وهو في
    الابتداء~$\OLMathNumeral{1}$.
    \item تقصد الكتلة رقم~$\OLId{latin-ordinary}{k}$ من «الشريط»، وتقرأ «الرمز» فيها.
    \item\ollabel{find-inst}%
    تطلب التعليمة الموافقة ل«الحالة» و«الرمز» الحاليين.
    \item تعود إلى الكتلة رقم~$\OLId{latin-ordinary}{k}$، وتجعل موضع «رمزها» عدد ترميز
    الرمز الذي تكتبه~$\OLId{latin-looped}{M}_\OLId{latin-ordinary}{e}$.
    \item تقصد موضع تسجيل «الحالة الحالية»، وتبدل عددها بعدد
    الحالة الجديدة.
    \item تقصد موضع تسجيل «موضع الشريط»، فتحذف علامة
    $\TMstroke$ أو تزيد علامة~$\TMstroke$، بحسب كون الحركة
    المأمور بها إلى اليسار أو إلى اليمين، على الترتيب.
    \item تعيد هذه الخطوات.\footnote{نترك في هذا الوصف بعض
    الدقائق؛ فقد تفتقر~$\OLId{latin-looped}{U}$ إلى موضع أوسع عند زيادة العداد الذي
    يسجل «موضع الرأس الحالي»، فيلزمها حينئذ نقل «الشريط» كله
    إلى اليمين.}
  \end{enumerate}
  فإن كانت $\OLId{latin-looped}{M}_\OLId{latin-ordinary}{e}$ لا تتوقف عند الدخل~$\OLId{latin-ordinary}{n}$، لم تتوقف $\OLId{latin-looped}{U}$ أيضًا،
  وكان خرجها غير معرّف.

  وأما إذا لم تجد في الخطوة~\olref{find-inst}، ضمن وصف~$\OLId{latin-looped}{M}_\OLId{latin-ordinary}{e}$،
  تعليمة لزوج «الحالة» و«الرمز» الحاليين، فذلك موضع توقف $\OLId{latin-looped}{M}_\OLId{latin-ordinary}{e}$.
  وعندئذ تمحو $\OLId{latin-looped}{U}$ ما يقع يسار «الشريط»، وتتخطى كتلته الأولى:
  علامة~$\TMstroke$ واحدة ترمز $\TMendtape$. ثم تبدأ بالكتلة
  التالية، فتضع في أقصى يسار شريطها علامة~$\TMstroke$ واحدة
  عن كل كتلة متعاقبة من ثلاث علامات~$\TMstroke$، وتمحو في أثناء
  ذلك ما فرغت منه من كتل «الشريط». فإذا بلغت كتلة من
  علامتي~$\TMstroke$، وهي ترميز $\TMblank$، فحصت كل ما بعدها من
  الخانات المرمّزة. فإن كانت كلها فراغات محت ما بقي وتوقفت؛ وإن
  وجدت بعد ذلك رمزًا غير فارغ تباعدت أو تركت خرجًا غير أحادي عمدًا،
  فلا تقبل ترميزًا غير قياسي. وبتمام العمل المقبول لا يبقى على شريط
  $\OLId{latin-looped}{U}$ إلا كتلة واحدة من علامات~$\TMstroke$ طولها~$\OLId{latin-ordinary}{m}$.

  % تصحيح مصدر (C267-01): أُضيف فحص ما بعد أول فراغ لرفض ترميزات الشريط غير القياسية.

  فإن وجدت $\OLId{latin-looped}{U}$ قبل بلوغ رمز~$\TMblank$ كتلة لا ترمز
  $\TMstroke$، أي ليست من ثلاث علامات~$\TMstroke$، توقفت من
  فورها. وحينئذ لا يكون ما على شريط~$\OLId{latin-looped}{U}$ كتلة واحدة من
  علامات~$\TMstroke$، فلا يكون خرجها عددًا طبيعيًا؛ أي إن
  $\OLId{latin-ordinary}{f}(\OLId{latin-ordinary}{e},\OLId{latin-ordinary}{n})$ غير معرّفة في هذه الحال.
\end{proof}

\end{document}
